% ARCHIVED at v3.66 (2026-09-22, author's instruction): every UAV-corridor passage of Chapters 5 and 6 as it stood at v3.65 (corridor_v2_slide corpus, batch_uav_20260919_111701 tag u17cv2 and the 18-09 rollout drop). The chapters now carry the SAME structure with the data blanked (pending) for the corridor-v3 rebuild (PENDING_20260922_corridor_v3_run.md).
% Verbatim copy of the v3.65 text, line-for-line, for reference only. NOT built. Do not edit.


% ===== Chapter 5, the UAV-corridor item of sec:setup:tasks:uav (05_setup.tex 333-341) =====
  \item[UAV-corridor] --- the entry test. Two parallel walls leave a clear width of $1.90$\,m. In the
    horizontal $x$--$y$ plane, a halfspace boundary slopes by $14^\circ$ from the upper wall at the
    entrance to $0.05$\,m below the centre line at the exit. It requires a lateral route change of
    $0.25$--$0.49$\,m before the exit. Every constraint of the scene, the walls included, is a vertical
    surface: it constrains the horizontal position only, and the altitude is bounded by nothing but the
    workspace box (\autoref{fig:uav-corridor-altitude}).
    The models are trained on the demonstrations of a narrower corridor and never see this halfspace, so it is
    a test-time constraint in the same sense as the halfspaces of D3IL-avoiding.
  \item[UAV-pillars] --- \flawed{temporarily disabled, v3.65: the scene is under review} six cylindrical pillars stand in two rows across three columns, and the vehicle

% ===== Chapter 5, tab:uav-scenes corridor rows (364-366) =====
    UAV-corridor
      & two wall halfspaces & inner faces $y=\pm0.95$\,m, active over $x\in[-2,2]$; clear width $1.90$\,m \\
      & sloping halfspace & from $(-2.0,\,0.95)$ to $(2.0,\,-0.05)$, a descent of $14^\circ$ \\

% ===== Chapter 5, sec:setup:data:uav corridor paragraph (563-567) =====
UAV-corridor is built the other way, and it is the way the obstacle-avoidance benchmark is built. The
three lanes clear both walls, but every one of them crosses the sloping halfspace, by $0.24$, $0.36$ and $0.47$\,m
for the left, centre and right lane --- because the halfspace is introduced at test time and did not exist
when the demonstrations were flown. The scene in which the constraint is new to the data is exactly the
scene in which the projection has something to do.

% ===== Chapter 5, tab:protocol corridor row (1167) =====
    UAV-corridor & 1 & 12 flights, four on each of the routes L, C and R & the twelve flights \\

% ===== Chapter 5, sec:setup:protocol:uav (1198-1216) =====
\subsubsection{Quadrotor Benchmark}\label{sec:setup:protocol:uav}

The three scenes are evaluated and reported separately, never pooled. UAV-corridor has twelve flights,
four on each of its left, centre and right routes, so each route contributes equally. UAV-corridor constrains the commanded
position in addition to the measured one and leaves the action unconstrained, so its projection configuration is not
the one used on the other two scenes. Every flight is tracked by the cascaded geometric controller
with the brake-to-rest velocity setpoint defined in \autoref{sec:method:geometric}; the one exception is the controller
comparison of \autoref{sec:res:uav:controller}, which replaces the controller by MuJoCo MPC and says so.
\srcnote{\texttt{Gen15/U16/DA\_20260914\_corridor\_v2\_paper\_full.md} \S1
(\texttt{-bounds\_free-pdes-tightened}, \texttt{pid\_stopgo}, $n=12$);
\texttt{Campaign\_20260907\_five\_missions/DA\_20260912\_pillars\_diffusion\_baseline\_reference.md} \S8;
\texttt{batch\_uav\_20260915\_100816/per\_rollout\_detail.csv}, column \texttt{controller}: all 960 rows on
\texttt{pillars\_hg} and 880 of the 889 rows on \texttt{s\_curve\_hg} are \texttt{pid\_stopgo}, the other 9
\texttt{mjpc}; \texttt{pid\_stopgo} sets \texttt{v\_des} $=0$ at each setpoint,
\texttt{eval\_mix\_uav.py:1598, 1716--1720}.}

% OPEN (author, 2026-09-20): this subsection may be dropped. Not decided. It is the worked example
% behind \eqref{eq:method:degen:ngen}; if it goes, the two settings read off below have to move into
% \autoref{sec:res:constraints:degenerate}, which is the only place that uses the count, and

% ===== Chapter 6, 6.3.1 UAV-corridor before projection, sec:res:uav:corridor:raw + tab:uav-corridor-raw (1341-1392) =====
\subsubsection{UAV-corridor}\label{sec:res:uav:corridor:raw}

Without projection, every model flies the route it learned and ignores the sloping halfspace, which it never
saw in training (\autoref{tab:uav-corridor-raw}). All four succeed on all twelve flights, and none of them
flies a single flight collision-free: $43$ to $56$ control steps per flight violate the halfspace, and the
count barely moves with the budget. The diffusion baseline is in that group --- unprojected it crosses the
finish line on every flight, in $264$ control steps, with $49.8$ violating steps per flight. Nothing here
orders the models, because the constraint fails by construction for every one of them; what the table fixes
is the price of generation alone, $9$ to $45$\,ms per control step for the flow-based models across the
ladder and $181$\,ms for the baseline at its budget, against which the projected costs of
\autoref{sec:res:uav:corridor} are read. No cost frontier is drawn for this arm: with success with
constraint satisfaction at $0.00$ in every cell there is nothing to trade against time.

\begin{table}[htpb]
  \caption[UAV-corridor: the generative models before projection]{UAV-corridor with the sloping test-time
  halfspace, \textbf{before projection}: the plan flown as generated, twelve flights per row (four per
  route), training seed 6. \emph{Success}: crosses the finish line. \emph{Collision-free}: no violating
  control step. \emph{S\&C}: success on a collision-free flight. \emph{Violating steps}: per flight,
  mean $\pm$ the sample standard deviation over the twelve, as is the time
  (\autoref{sec:setup:metrics:spread}); the three count columns are counts. \baselinemark\ the diffusion
  model of \ac{DPCC}. The same cells after projection are \autoref{tab:uav-corridor}.}
  \label{tab:uav-corridor-raw}
  \centering
  \footnotesize
  \setlength{\tabcolsep}{3pt}
  \begin{tabularx}{\linewidth}{@{}L r rrr rr@{}}
    \toprule
    Model & $\nfe$ & Success & \shortstack{Collision-\\free} & S\&C & \shortstack{Violating\\steps} & ms/step \\
    \midrule
    MeanFM & 1 & 12/12 & 0/12 & 0/12 & $54.2 \pm 14.3$ & $9.5 \pm 0.3$ \\
    MeanFM & 3 & 12/12 & 0/12 & 0/12 & $54.8 \pm 12.7$ & $27.2 \pm 0.3$ \\
    MeanFM & 5 & 12/12 & 0/12 & 0/12 & $54.6 \pm 12.6$ & $45.2 \pm 0.4$ \\
    \midrule
    CI-MeanFM, $\alpha_{\mathrm{end}}=0.2$ & 1 & 12/12 & 0/12 & 0/12 & $54.8 \pm 12.6$ & $9.4 \pm 0.3$ \\
    CI-MeanFM, $\alpha_{\mathrm{end}}=0.2$ & 3 & 12/12 & 0/12 & 0/12 & $56.1 \pm 12.3$ & $26.9 \pm 0.4$ \\
    CI-MeanFM, $\alpha_{\mathrm{end}}=0.2$ & 5 & 12/12 & 0/12 & 0/12 & $56.2 \pm 12.3$ & $44.6 \pm 0.3$ \\
    \midrule
    FM & 1 & 12/12 & 0/12 & 0/12 & $42.6 \pm 12.1$ & $9.1 \pm 0.3$ \\
    FM & 3 & 12/12 & 0/12 & 0/12 & $46.1 \pm 12.7$ & $26.7 \pm 0.9$ \\
    FM & 5 & 12/12 & 0/12 & 0/12 & $47.6 \pm 12.8$ & $43.7 \pm 0.5$ \\
    \midrule
    \mainconfig{\baselinemark\ Diffusion} & \mainconfig{20} & \mainconfig{12/12} & \mainconfig{0/12} & \mainconfig{0/12} & \mainconfig{$49.8 \pm 13.1$} & \mainconfig{$180.6 \pm 1.5$} \\
    \bottomrule
  \end{tabularx}
\end{table}
\dataref{analysis\_results\_checkpoint/19-09-UAV-Pillars-Exclude/batch\_uav\_20260919\_111701, tag u17cv2,
variant \texttt{diffuser}, geometry \texttt{corridor\_cv2s}: \texttt{success\_relaxed},
\texttt{collision\_free\_completed}, \texttt{n\_success\_relaxed\_and\_constraints}, \texttt{n\_violations},
\texttt{avg\_time\_ms}; recounted 2026-09-22 (v3.63).}

\FloatBarrier
\begin{pillarsflawed}

% ===== Chapter 6, 6.3.2 UAV-corridor after projection, sec:res:uav:corridor: tab:uav-corridor, fig:uav-corridor-tradeoff, fig:uav-corridor-paths, Altitude paragraph + fig:uav-corridor-altitude (1461-1674) =====
\subsubsection{UAV-corridor}\label{sec:res:uav:corridor}

With per-step projection the models separate, and they separate on the constraint rather than on the goal
(\autoref{tab:uav-corridor}; the same cells before projection are \autoref{tab:uav-corridor-raw}). Analytic average-velocity matching and consistency-interpolated
average-velocity matching succeed collision-free on all twelve flights at $\nfe=3$ and at $\nfe=5$.
Instantaneous-velocity matching succeeds on every flight too, but not collision-free: it keeps $1.5$
to $5.2$ shallow violating steps per flight and is collision-free on at most two of twelve. The diffusion
baseline is collision-free on every flight and succeeds on none.

That last row is the one to read carefully, because the flown paths do not look like a failure
(\autoref{fig:uav-corridor-paths}). The finish line is not a vertical plane at the end of the corridor: it
is the plane through the end of the demonstrated route, perpendicular to that route's final approach
heading, and the demonstrated routes turn downward as they leave the corridor. Projected, the diffusion
baseline holds the lane it learned and never makes that turn. It therefore travels as far along $x$ as
the models that succeed --- its flights end between $x=2.66$ and $x=2.79$\,m, inside the range of the
successful flights --- while staying on the near side of the line, and it runs the full $396$-step episode limit on
every flight rather than stopping early on arrival, ending $0.51$\,m from the goal point. Collision-free
and $0/12$ is not a contradiction here: the projector kept it out of the halfspace, and it did so by
keeping it on a lane that does not reach the goal.
\srcnote{\texttt{success\_relaxed = crossed\_line and safe}
(\texttt{mix\_uav\_test/eval\_mix\_uav.py:1872}); the line is built at \texttt{:1580--1588} through the
expert endpoint, normal to the heading over the last $0.1$\,s of the reference, and latches on
\texttt{:1757}. \texttt{safe} is contact-free and airborne (\texttt{:1849}) and is $1$ for every diffusion
flight here, so the failure is the crossing test alone, not a crash. Corpus values for the cell:
\texttt{goal\_crossed\_line} $=0$, \texttt{phys\_safe} $=1$, \texttt{constraint\_collision\_free} $=1$,
\texttt{n\_steps} $=396=$ \texttt{max\_episode\_length}, \texttt{goal\_dist} $=0.5099$; the unprojected arm
of the same cell has \texttt{goal\_crossed\_line} $=1$ and $49.8$ violating steps.}

No smaller budget of the diffusion baseline was evaluated on this scene. Its step count is fixed when its
noise schedule is discretised at training, so each additional budget is a separate training run
(\autoref{sec:res:avoiding:models}), and at the budget it was trained for it already fails the scene's
goal criterion on every flight. A cheaper checkpoint of the same model cannot recover a goal the expensive
one does not reach, so the runs were not spent. Analytic average-velocity matching at $\nfe=3$ is therefore ahead of instantaneous-velocity
matching on the constraint, $12/12$ against $0/12$, and ahead of the diffusion baseline on the goal,
$12/12$ against $0/12$, at $0.16$ of its time per control step.

\begin{table}[htpb]
  \caption[UAV-corridor: the generative models after projection]{UAV-corridor with the sloping test-time
  halfspace, \textbf{after projection}: every row is per-step projection with temporal consistency on
  tightened constraints, the configuration every later result of this scene uses; twelve flights per row
  (four per route). The same cells before projection are \autoref{tab:uav-corridor-raw}. \emph{Success}:
  crosses the finish line. \emph{Collision-free}: no violating control step. \emph{S\&C}: success on a
  collision-free flight. \emph{Violating steps}: per flight, mean $\pm$ the sample standard deviation over
  the twelve, as is the time (\autoref{sec:setup:metrics:spread}); the three count columns are counts.
  Training seed 6. \selectedmark\ the selected configuration, \baselinemark\ the diffusion model of
  \ac{DPCC}.}
  \label{tab:uav-corridor}
  \centering
  \footnotesize
  \setlength{\tabcolsep}{3pt}
  \begin{tabularx}{\linewidth}{@{}L r rrr rr@{}}
    \toprule
    Model & $\nfe$ & Success & \shortstack{Collision-\\free} & S\&C & \shortstack{Violating\\steps} & ms/step \\
    \midrule
    MeanFM                         & 1  & 12/12 & 2/12  & 2/12  & $1.8 \pm 1.1$ & $29.1 \pm 0.5$ \\
    \mainconfig{\selectedmark\ MeanFM} & \mainconfig{3} & \mainconfig{12/12} & \mainconfig{12/12} & \mainconfig{12/12} & \mainconfig{$0.0 \pm 0.0$} & \mainconfig{$99.1 \pm 1.1$} \\
    MeanFM                         & 5  & 12/12 & 12/12 & 12/12 & $0.0 \pm 0.0$ & $143.9 \pm 3.7$ \\
    \midrule
    CI-MeanFM, $\alpha_{\mathrm{end}}=0.2$ & 1  & 12/12 & 4/12  & 4/12  & $1.2 \pm 1.2$ & $29.0 \pm 0.3$ \\
    CI-MeanFM, $\alpha_{\mathrm{end}}=0.2$ & 3  & 12/12 & 12/12 & 12/12 & $0.0 \pm 0.0$ & $101.9 \pm 4.1$ \\
    CI-MeanFM, $\alpha_{\mathrm{end}}=0.2$ & 5  & 12/12 & 12/12 & 12/12 & $0.0 \pm 0.0$ & $141.8 \pm 1.7$ \\
    \midrule
    FM                             & 1  & 12/12 & 0/12  & 0/12  & $5.2 \pm 0.8$ & $29.2 \pm 0.6$ \\
    FM                             & 3  & 12/12 & 0/12  & 0/12  & $2.1 \pm 0.8$ & $86.4 \pm 1.5$ \\
    FM                             & 5  & 12/12 & 2/12  & 2/12  & $1.5 \pm 1.0$ & $133.5 \pm 4.7$ \\
    \midrule
    \mainconfig{\baselinemark\ Diffusion} & \mainconfig{20} & \mainconfig{0/12} & \mainconfig{12/12} & \mainconfig{0/12} & \mainconfig{$0.0 \pm 0.0$} & \mainconfig{$625.5 \pm 22.0$} \\
    \bottomrule
  \end{tabularx}
\end{table}

\dataref{analysis\_results\_checkpoint/19-09-UAV-Pillars-Exclude/batch\_uav\_20260919\_111701, tag u17cv2, recomputed by
DA\_in\_Paper/analysis/uav\_results.py; agrees with Gen15/U16/DA\_20260914\_corridor\_v2\_paper\_full \S3--\S7.
The unprojected twins of these cells are \autoref{tab:uav-corridor-raw}.}
\guard{The three flow-based projected rows are the projection variant \texttt{dpcc-t-tightened}. The diffusion row
is \texttt{dpcc-t-bounds\_free-pdes-tightened}: the same per-step projector and the same tightened
geometry, with the action-magnitude bound released and the plan anchored on the commanded rather than the
measured position. It is the only projected corridor variant the baseline was evaluated under, so the
comparison holds the projector and the constraint set fixed but not the two auxiliary switches.
Verified against the corpus 2026-09-21.}

The budgets are $\nfe\in\{1,3,5\}$ rather than $\{1,2,5\}$ because of what the second one has to support.
At the activation threshold of $0.5$ this scene runs, the projector acts on the last
$\lceil \eta\nfe \rceil$ sampling steps, and at $\nfe=1$ and $\nfe=2$ that is the terminal step alone ---
the degenerate case of \autoref{sec:res:constraints:degenerate}, in which endpoint projection is
projection after sampling and has nothing to guide. $\nfe=3$ is the smallest budget at which it has a
real guiding step, so it is the smallest budget at which the two projectors of
\autoref{sec:res:uav:projection} can be compared at all. Choosing it for the model table as well keeps
one budget ladder across the whole scene instead of two.

The same numbers against their price give this scene its cost frontier
(\autoref{fig:uav-corridor-tradeoff}). As on D3IL-aligning, and unlike D3IL-avoiding, the control-step
count is not the cost to plot: a flight that succeeds keeps flying to the episode limit, so its
step count measures the limit and not the arrival. What is plotted is therefore success with constraint
satisfaction against the time one action costs.

All three selection rules of every cell are tabled in \autoref{tab:uav-corridor-projection}, and points
that would print on top of each other are displaced by a few pixels, which moves the drawing and not the
data. The success axis is logarithmic in the number of successful flights, so that the difference between
two and six of twelve is as visible as the difference between six and twelve; a configuration that never
succeeds cannot be drawn on it and sits on a floor below the axis break. Two rules follow from that
drawing: a configuration with no successful flight is not a trade-off and is not eligible for the
frontier, and two costs within one per cent of each other are treated as one cost, so that a tenth of a
millisecond cannot put a point on the frontier.

Four configurations are on the frontier and they are the two average-velocity models at $\nfe=3$ and at
$\nfe=5$, all four at $1.00$. The cheapest of them costs $97.7$\,ms per control step, and the fifth point
on the staircase is analytic average-velocity matching at $\nfe=1$, which buys $0.50$ for $29.2$\,ms ---
so a third of the cost carries half the flights and no more. Nothing else is on it: instantaneous-velocity
matching reaches at most $0.17$ at any budget, and the diffusion baseline is the rightmost point in the
figure at $625$\,ms and $0.00$, because within the episode it never reaches the finish line. The shape of
this frontier is a step rather than a slope, which is the difference between this scene and
D3IL-aligning: here the budget either buys the constraint or does not.

\begin{figure}[htpb]
  \centering
  \includegraphics[width=0.92\linewidth]{fig_uav_corridor_tradeoff}
  \caption[UAV-corridor: outcome against cost]{UAV-corridor, success with constraint satisfaction, as
  successful flights of twelve on a log axis, against the time to compute one action, on a log axis;
  per-step projection, tightened constraints, training seed 6, twelve
  flights per configuration. Configurations with no successful flight are drawn on the floor below the
  axis break. Each model is drawn once per step budget, at the selection rule that scores
  highest there, and marker shape says which rule that was. A ring marks a configuration with at least one
  successful flight that nothing else beats on both axes, costs within one per cent counting as equal, and
  the dashed staircase joins them; hollow markers are dominated or never succeed, the number beside
  a point is its step budget, and the boxed arrow gives the direction in which both axes improve.
  Legend: MeanFM, CI-MeanFM, FM and the diffusion model of \ac{DPCC}.}
  \label{fig:uav-corridor-tradeoff}
\end{figure}
\dataref{analysis\_results\_checkpoint/19-09-UAV-Pillars-Exclude/batch\_uav\_20260919\_111701, tag
u17cv2, variants \texttt{dpcc-\{r,c,t\}-bounds\_free-pdes-tightened}; built by
DA\_in\_Paper/plotting/builders/frontier.py. Every point reproduces its cell of
\autoref{tab:uav-corridor-projection}.}

UAV-pillars (\flawed{disabled, v3.65}) is not drawn as a frontier: \autoref{tab:uav-pillars-best} already places each model at its best
configuration per projection block and budget, and the scene supports no single selected configuration
(\autoref{sec:res:uav:pillars:projected}). UAV-s-curve has
data, but every configuration it has --- projected and unprojected, at every budget flown --- scores
$0.00$ success with constraint satisfaction (\autoref{tab:uav-scurve}, \autoref{tab:uav-scurve-projection}).
A cost frontier over points that all share one outcome is a ranking of failures by price, so what that
scene gets instead is the abort count printed beside each of its rows.

The budget decides how much of the halfspace violation the projection removes. At $\nfe=1$ a single projected sampling
step reduces the violations but does not remove them: analytic average-velocity matching is collision-free
on two flights of twelve, consistency-interpolated average-velocity matching on four, and
instantaneous-velocity matching on none. From $\nfe=3$ on, both average-velocity models fly every flight
clean. UAV-corridor thus establishes that the flow-based planners fly a constrained course that the
diffusion baseline, under the same projection, does not finish.
Every row above is per-step projection, because on this scene that is the method:
\autoref{sec:res:uav:projection} runs both projectors at $\nfe=3$ and $\nfe=5$ under every selection
rule, and per-step projection reaches $1.00$ in all twelve cells of the two average-velocity models
where endpoint projection reaches $0.58$ at best (\autoref{tab:uav-corridor-projection}). The choice of
projector is therefore settled there and held fixed here, so that this subsection compares generative
models and nothing else.

\autoref{fig:uav-corridor-paths} draws those flights. A flight can succeed through an obstacle, so
its two colour codes are read separately. It shows the trade the table reports as two numbers:
the budget at which a model stops crossing the sloping halfspace is also the budget at which it stops reaching the end
of the corridor, and the diffusion baseline sits at the far end of that trade, flying twelve clean flights
that all stop short.

\begin{figure}[htpb]
  \centering
  \includegraphics[width=\textwidth]{fig_uav_corridor_paths}
  \caption[UAV-corridor: flown paths per model and budget]{UAV-corridor seen from above, drawn like
  \autoref{fig:constraints-uav}, with the executed path of each of the twelve flights over it. One row
  per model across the budgets $\nfe\in\{1,3,5\}$, the diffusion baseline last at its own budget of
  twenty; every panel is the configuration \autoref{tab:uav-corridor} reports. Colour reports the
  constraints (green: collision-free; red: entered an inflated obstacle), and the stroke and end mark
  report success (solid and filled: success; dotted and hollow: no success). Panel titles: MeanFM,
  CI-MeanFM, FM, Diffusion.}
  \label{fig:uav-corridor-paths}
\end{figure}
\dataref{DA\_in\_Paper/data/uav\_paths.json, extract/uav\_paths.py; ten cells of
batch\_uav\_20260915\_100816, tag u17cv2}

\paragraph{Altitude.} Every constraint of this scene acts on the horizontal position. The two walls, the
sloping halfspace and the wall-end caps are vertical surfaces, and the altitude of the plan is bounded by
nothing but the workspace box, whose limits after inflation --- $0.61$ and $1.49$\,m --- lie $0.3$\,m or
more from any altitude that was flown. The corridor therefore tests a lateral route change and nothing in
$z$. \autoref{fig:uav-corridor-altitude} shows the flights of \autoref{fig:uav-corridor-paths} from the
side, without and with projection, so that this is seen rather than assumed. The projection does move the
altitude: launched from the same $1.11$\,m, the unprojected flights of the two average-velocity models sink
by $0.08$\,m over the corridor while their projected flights hold $1.10$--$1.11$\,m, and flight by flight
the two arms differ by up to $0.14$--$0.22$\,m for the flow-based models and $0.27$\,m for the diffusion
baseline, the difference growing along $x$ towards the exit, where the slide is tightest. That is a
by-product of the lateral correction --- the sample is regenerated after each correction, and the
configuration of this scene releases the action box (\autoref{sec:setup:protocol:uav}), so $z$ is free to
follow --- and not a response to a constraint, because none acts on $z$ within the flown band.
\guard{Whether the projector clears a constraint that acts on altitude is not tested by this scene; a
corridor with a sloping surface in $x$--$z$ is planned (Gen15 U19), and no claim here rests on it.}

\begin{figure}[htpb]
  \centering
  \includegraphics[width=\textwidth]{fig_uav_corridor_altitude}
  \caption[UAV-corridor: flown altitude, without and with projection]{UAV-corridor seen from the side:
  the altitude $z$ of every flight over its position $x$ along the corridor, one panel per model at
  $\nfe=3$ and the diffusion baseline at $\nfe=20$; grey without projection, blue under the per-step
  projection of \autoref{tab:uav-corridor}; twelve flights per arm, training seed 6. The shaded span is
  the length of the walls and of the sloping halfspace, both of which act on the horizontal position only.
  The dashed lines are the only altitude limits the projector sees, the workspace box after inflation by
  the rotor reach. The altitude axis is stretched: $1.2$\,m of altitude spans the height of a panel that
  is $5.6$\,m long. A dot marks where a flight ended, filled for success. Panel titles: MeanFM, FM,
  CI-MeanFM, Diffusion.}
  \label{fig:uav-corridor-altitude}
\end{figure}
\dataref{DA\_in\_Paper/data/uav\_paths.json, group \texttt{corridor\_altitude}, extract/uav\_paths.py
(\texttt{obs\_all} dim 5); eight cells of batch\_uav\_20260915\_100816, tag u17cv2, variants
\texttt{diffuser} and \texttt{dpcc-t-bounds\_free-pdes-tightened}; box limits
config/uav\_projection.yaml \texttt{corridor\_v2\_slide.workspace\_bounds}; the sink, hold and paired
differences recounted from the same rollouts 2026-09-22 (v3.65).}

\FloatBarrier
\begin{pillarsflawed}

% ===== Chapter 6, 6.3.3 Projection Methods, the corridor part: prose + tab:uav-corridor-projection + tightening result (2024-2096) =====
\subsection{Projection Methods}\label{sec:res:uav:projection}

On UAV-corridor, per-step projection is the method, and the comparison is now complete: both methods have
been run at $\nfe=3$ and $\nfe=5$ under every selection rule, for all three flow-based models
(\autoref{tab:uav-corridor-projection}). Per-step projection succeeds collision-free on all twelve
flights for both average-velocity models under all three rules at both budgets. Endpoint projection reaches
$0.58$ at best --- analytic average-velocity matching at $\nfe=3$ with cumulative projection cost --- and
is below half on every other cell of those two models. It is the cheaper arm, $74$\,ms against $99$\,ms per
control step at $\nfe=3$ and $47$\,ms in its single-candidate form, so what it offers on this scene is
about half the projection cost for about half the constraint satisfaction. The bar this thesis set for it
--- overtaking the per-step projector of \ac{DPCC} at a lower projection threshold --- is not cleared here.

Endpoint projection also gets worse as the budget grows, where per-step projection does not move: analytic
average-velocity matching falls from $0.58$ to $0.33$ and consistency-interpolated from $0.42$ to $0.17$
between $\nfe=3$ and $\nfe=5$, while per-step projection stays at $1.00$ throughout. Raising the budget
lengthens the ODE whose endpoint the single solve has to anticipate, so the prediction it corrects is
further from what the sampler finally produces.

\begin{table}[htpb]
  \caption[UAV-corridor: the two projection methods under every selection rule]{UAV-corridor, success with
  constraint satisfaction --- success on a collision-free flight --- for both projection methods
  under every selection rule, twelve flights per cell, tightened constraints, training seed 6.
  \emph{Per-step}: the projector of \ac{DPCC} under random selection ($r$), cumulative projection cost
  ($c$) and temporal consistency ($t$). \emph{Endpoint}: endpoint projection, in its single-candidate form
  and under the same three rules. A dash is a configuration that was not evaluated.
  \selectedmark\ the selected configuration of this scene.}
  \label{tab:uav-corridor-projection}
  \centering
  \footnotesize
  \setlength{\tabcolsep}{5pt}
  \begin{tabularx}{\linewidth}{@{}L r rrr rrrr@{}}
    \toprule
    & & \multicolumn{3}{c}{Per-step} & \multicolumn{4}{c}{Endpoint} \\
    \cmidrule(lr){3-5}\cmidrule(l){6-9}
    Model & $\nfe$ & $r$ & $c$ & $t$ & single & $r$ & $c$ & $t$ \\
    \midrule
    \mainconfig{MeanFM} & \mainconfig{3} & \mainconfig{1.00} & \mainconfig{1.00} & \mainconfig{\selectedmark\ 1.00} & \mainconfig{0.42} & \mainconfig{0.50} & \mainconfig{0.58} & \mainconfig{0.42}  \\
           & 5 & 1.00 & 1.00 & 1.00 & 0.00 & 0.08 & 0.33 & 0.17 \\
    \midrule
    CI-MeanFM, $\alpha_{\mathrm{end}}=0.2$ & 3 & 1.00 & 1.00 & 1.00 & 0.25 & 0.33 & 0.42 & 0.58 \\
           & 5 & 1.00 & 1.00 & 1.00 & 0.00 & 0.08 & 0.00 & 0.17 \\
    \midrule
    FM     & 3 & 0.00 & 0.00 & 0.00 & 0.00 & 0.00 & 0.00 & 0.00 \\
           & 5 & 0.00 & 0.08 & 0.17 & 0.00 & 0.00 & 0.00 & 0.00 \\
    \midrule
    \mainconfig{\baselinemark\ Diffusion} & \mainconfig{20} & \mainconfig{0.00} & \mainconfig{0.00} & \mainconfig{0.00} & --- & --- & --- & ---  \\
    \bottomrule
  \end{tabularx}
\end{table}
\dataref{analysis\_results\_checkpoint/19-09-UAV-Pillars-Exclude/batch\_uav\_20260919\_111701, tag u17cv2,
variants \texttt{dpcc-\{r,c,t\}} and \texttt{hardflow\_sls-\{r,c,t\}}; uav\_results.py; DA\_20260919 \S2.
The \texttt{hardflow\_sls-r} and \texttt{hardflow\_sls-c} cells are new in that batch.}
\guard{Endpoint projection has no row at $\nfe=1$ on this scene and will not get one: at a single ODE step
the only step is the terminal one, so there is no guiding step for the method to act on
(\autoref{sec:res:constraints:degenerate}). That is a property of the method at that budget, not a pending
evaluation. The baseline's four dashes are the opposite case: endpoint projection has never been run on
the diffusion model on any scene, so the method is compared against the per-step projector on the
flow-based models alone. Its per-step row is printed because the scene's two phases should cover the same
models, and it says what \autoref{tab:uav-corridor} already says --- the baseline satisfies the
constraints and does not reach the finish line.}

The same table carries a result about the projector rather than the method. The untightened companion of
the per-step configuration scores $0.00$ at every budget for every model, against $1.00$ for its tightened
twin, with $48.1$ violating control steps per flight for analytic average-velocity matching at $\nfe=3$. On
this scene the tightening margin of \ac{DPCC} is not a refinement of the projector; it is the whole of its
constraint satisfaction, and every corridor result of this chapter is reported on tightened constraints for
that reason.
\dataref{same batch, variant \texttt{dpcc-c-bounds\_free-pdes}; DA\_20260919 \S2}

The projection configuration of this scene differs from that of the other quadrotor scene
(\autoref{sec:setup:protocol:uav}), so the scenes are not pooled.

\begin{pillarsflawed}

% ===== Chapter 6, 6.3.5 UAV conclusion (2288-2318) =====
\subsection{Conclusion}\label{sec:res:uav:conclusion}

UAV-corridor shows that the flow-based planners fly under constraints: analytic and
consistency-interpolated average-velocity matching with per-step projection at $\nfe=3$ succeed
collision-free on every flight, at a sixth of the time of the diffusion baseline, which stays collision-free
but never reaches the finish line. \flawed{The UAV-pillars sentence that follows is disabled, v3.65.}
UAV-pillars, on which the plans are feasible before projection, ranks
nothing among the three flow-based models: instantaneous-velocity matching with endpoint projection succeeds
collision-free on all ten flights at $78$\,ms per control step, analytic average-velocity matching on nine
of ten in a shorter path, and at $\nfe=2$ the order is reversed; the one ordering that holds in every cell
is against the diffusion baseline, $0.30$ at best at $4.4$\,s per control step. UAV-s-curve has no
supported combination; its limit is the tracking controller, and the scene's flights end with the vehicle
inverted rather than off course. What the
quadrotor benchmark supports is therefore UAV-corridor's result, and it is a strong one: a test-time constraint the demonstrations do not satisfy,
under which both average-velocity models at $\nfe=3$ succeed collision-free on every flight, and the
diffusion baseline on none.

% =============================================================================
% =============================================================================
% 🔒 DO NOT EDIT (author, 2026-09-21 / v3.60)
%
% This section summarises 6.1, 6.2 and 6.3. It is the LAST thing to be written,
% because every number in it is a restatement of a number from those three, and
% a summary written over sections that are still moving has to be rewritten each
% time one of them moves.
%
% Leave it exactly as it stands until 6.1, 6.2 and 6.3 have each been verified
% against the corpus and locked by the author. Then reconcile this section with
% them in one pass. Do not "fix" a cell here on its own -- if a number here
% disagrees with its source section, the source section is what to check.
% =============================================================================
